\documentclass[11pt]{article}
\usepackage[margin=1in]{geometry}
\usepackage{amsmath,amssymb,amsthm}
\usepackage{booktabs,array,longtable}
\usepackage{hyperref}
\usepackage{url}
\usepackage{microtype}
\usepackage{xcolor}
\hypersetup{colorlinks,linkcolor=black,citecolor=blue!50!black,urlcolor=blue!50!black}
\newcommand{\F}{\mathbb{F}}
\newcommand{\Z}{\mathbb{Z}}
\newcommand{\Q}{\mathbb{Q}}
\newcommand{\End}{\operatorname{End}}
\newcommand{\ord}{\operatorname{ord}}
\newcommand{\lrand}{\texttt{lrand48}}
\newcommand{\todo}[1]{\textcolor{red}{[#1]}}

\title{The Skipped Candidates in the NIST and ANSI X9.62 Curve Seed Runs:\\ A Negative Result, and a Structural Observation on the K-163 Generator}
\author{Dan Shumow\thanks{\todo{affiliation}; \texttt{shumow@gmail.com}. Computations and drafting assisted by Claude (Anthropic).}}
\date{October 2026 --- draft}

\begin{document}
\maketitle

\begin{abstract}
In October 2026 Anthropic recovered the NIST and ANSI prime-curve seeds as C-library \lrand{} output seeded by Unix timestamps. The reproduction exposes a loose end: in four prime-field runs, eleven unpublished candidates give prime-order curves of the right size, while eight were published. We regenerate the runs, check the group orders, and compare roughly fifty properties of the nineteen viable curves (including, after review, the complete factorization of $b$ as an integer and deep ECM on the cubic, quartic and sextic extension-order cofactors), together with models of the generating program and a set-cover search for unions of small defects. No tested rule separates the groups. An exact comparison preserving the published counts and the final published candidate in each run admits a cover as small as the observed one in 88 of 126 relabellings. These exploratory negative results do not exclude untested numerical criteria, run-specific code, or variable RNG consumption. Seven published coefficients match the textbook square-root convention; recovered seed offsets constrain cumulative RNG alignment, not every call. Separately, the subgroup half of the K-163 generator has a 126-bit $x$-coordinate on a 163-bit curve. A coefficient-order error in the original decimal export has been corrected without changing that bit length. The short coordinate survives independent Python and PARI checks and remains a concrete provenance lead, but its source and construction are unidentified.
\end{abstract}

\section{Introduction}

ECDSA's standard curves were handed to NIST in 1997--98 with seed values and no account of their origin \cite{fips186-2,x962}. Koblitz and Menezes catalogued the resulting unease \cite{riddle}; Bernstein, Lange and coauthors showed how much freedom a seed-search would have given a malicious designer \cite{bada55}; Valsorda posted a bounty for the seeds' derivation \cite{bounty}. In October 2026 Anthropic's Frontier Red Team reported that the seeds are \lrand{} output: five NIST prime-curve seeds and five ANSI X9.62 example-curve seeds are 20 consecutive bytes of $\lrand() \gg 23$ after $\texttt{srand48}(t)$ with $t$ a Unix timestamp between May 1997 and April 1998, and four of the five binary-curve seeds are five 31-bit \lrand{} words from one sitting on 17 April 1998 \cite{anthropic}. The B-163 seed remains unexplained. Valsorda reproduced the result; the Anthropic team claimed the bounty and directed it to ISRG \cite{seeds-solved}. The work is one output of a broader programme of model-assisted cryptanalysis at Anthropic \cite{anthropic-crypto}, whose authors identify human validation of model-found results as the current bottleneck; the present paper is, among other things, such a validation. The same Anthropic post \cite{anthropic} also explained six of the base points: for P-192 and the five ANSI X9.62 example curves prime192v2, prime192v3, prime239v1, v2 and v3, the generator's $x$-coordinate is output of the \texttt{zrandom} generator in Lenstra's FreeLIP library \cite{freelip} (state $\times 3^{121} \bmod 2^{107}-1$), seeded with a small integer, with the first on-curve $x$ taken; those six were generated in May 1997. We rely on that finding below (Sections~\ref{sec:program} and~\ref{sec:k163}) and claim no part of it. The base points of P-224, P-256, P-384, P-521 and all ten binary curves were left unexplained there.

Both write-ups note a residue. A run produces candidates in order; a candidate is \emph{valid} if $-27/c$ is a square mod $p$ (so a coefficient $b$ exists), \emph{viable} if the resulting curve has prime order $n$ with $\lceil\log_2 n\rceil \le \lceil\log_2 p\rceil$ (an undocumented but evidently applied rule \cite{seeds-solved}). In the P-256 run, candidate 876 is viable but candidate 1842 was published; the P-224 run passed over three viable candidates and the two ANSI runs seven more. Valsorda summarizes: ``roughly twice as many valid seeds as were published,'' with a second undocumented criterion that ``still escapes us.''

This paper attempts to find that criterion. The tested properties and program models do not provide a separator. This does not exhaust cheaply computable properties or establish statistical equivalence. Section~\ref{sec:setup} fixes the data. Section~\ref{sec:invariants} reports the invariant sweep, Section~\ref{sec:program} the program models, Section~\ref{sec:cover} the set-cover test. Section~\ref{sec:k163} reports the K-163 observation. Section~\ref{sec:open} lists what remains open.

\section{Setup and reproduction}\label{sec:setup}

\paragraph{Runs.} The four prime-field runs and their published candidates are: P-256 (\texttt{srand48} 1997-08-16 01:52:34 UTC, candidate 1842); P-224 (1998-04-03 12:52:07, candidate 7026); the prime192 run (1997-05-18 14:45:34; P-192 at 139, prime192v3 at 5228, prime192v2 at 6013); the prime239 run (1997-05-15 15:49:03; prime239v1/v2/v3 at 2185, 2320, 3048).

\paragraph{Reproduction.} For every candidate we derive $c$ by the X9.62 SHA-1 expansion and both square roots of $-27/c$. Candidates with 2-torsion (the cubic has a root) are discarded without counting. Group orders were computed with full SEA in PARI/GP~2.10 built from source and checked by $n\cdot P = O$ on random points. For prime $n>4\sqrt p$ in the Hasse interval, one nonzero point killed by $n$ determines $\#E$; for composite orders these are consistency checks, not independent certificates. (An early-abort pass produced trace-sign errors in that build and was discarded.) Review checks with PARI/GP~2.19.0 certify all nineteen viable prime orders by primality, the Hasse bound, and scalar multiplication.

\paragraph{Square-root convention.} For every published curve over $p\equiv 3\pmod 4$---P-256, P-192, prime192v2, prime192v3, prime239v1, v2, v3---the published $b$ equals $(-27/c)^{(p+1)/4} \bmod p$, never the other root. This is consistent with the textbook routine, not proof that the other root was never tested. Accepting either root adds $4+10+9$ candidates across those runs. For P-224 ($p\equiv1\pmod4$) the roots give isomorphic curves but different integer coefficients and coordinates. Generator root 1 is the negative of the published P-224 coefficient. The coefficient study uses root 2 as its publication-aligned model and retains root-1 witnesses for all four candidates; the historical skipped-candidate convention is unknown. Published-generator scans use the exact standard tuple.

\paragraph{The nineteen.} With that convention and the size rule, the viable candidates are
\begin{center}
\begin{tabular}{llll}
\toprule
run & viable & published & skipped \\
\midrule
P-256 & 876, 1842 & 1842 & 876 \\
P-224 & 777, 1485, 5736, 7026 & 7026 & 777, 1485, 5736 \\
prime192 & 139, 2198, 5228, 5405, 6013 & 139, 5228, 6013 & 2198, 5405 \\
prime239 & 36, 1562, 2185, 2320, 2491, 2682, 2863, 3048 & 2185, 2320, 3048 & 36, 1562, 2491, 2682, 2863 \\
\bottomrule
\end{tabular}
\end{center}
eight published, eleven skipped---exactly the lists in \cite{anthropic,seeds-solved}. In every run the published candidate with the largest index is the last viable candidate of the run; this is where runs stopped, not a selection rule, since 139 and 5228 (and 2185, 2320) were published with later viable candidates skipped.

\paragraph{Calibration.} The unknown criterion rejected about half of the viable candidates, not one in a thousand of all candidates. Any proposed rule must fail at that rate.

\section{Curve invariants}\label{sec:invariants}

Every property below was computed for all nineteen curves and compared between the published and skipped groups. The strongest form of the negative evidence is in the two runs that published three curves each: there the published curves disagree \emph{with one another} on most binary properties, so no such property can have been required. Tables are in Appendix~\ref{app:tables}; the full record with scripts is in the repository.

\subsection{Arithmetic of $n$ and $t$}
Twist order $n' = 2p+2-n$: small-factor structure and cofactor primality interleave (published P-224: $3^2\cdot11\cdot47\times$ 212-bit composite; skipped 1485: $3\cdot31\times$ 218-bit composite). Complete factorizations of $n\pm1$ (PARI MPQS; all cofactors fell in seconds): the published P-256 has the smoothest $n-1$ of its run (largest prime factor 92 bits) while the candidate it displaced has $n-1 = 2\cdot3\cdot7\cdot11\cdot43\times P_{242}$; in prime239 the published largest prime factors (69, 128, 110 bits) sit inside the skipped range (76--219). Congruences of $t,n,b$ modulo 2, 3, 4, 5, 8, 16; $b<p/2$; $(n-1)/2$ prime: nothing.

\subsection{CM structure}
With $D = t^2-4p = f s^2$, $f$ fundamental, the conductor $s$ of $\Z[\pi]$ in $\mathcal{O}_K$ is $1$ for sixteen curves, $3$ for the published P-224, and $5$ for one published (prime239v2) and one skipped (2682) prime239 curve. The class number, estimated as $\sqrt{|f|}\,L(1,\chi_f)/\pi$ with the Euler product to $10^6$, is \emph{smallest} in its run for the published P-224 and \emph{larger} for the published P-256 than for its skipped competitor; in the ANSI runs the estimates interleave. No tested universal floor or ceiling separates the groups. These estimates have no certified truncation error; stability at two Euler-product cutoffs is not an error bound. The heuristic $\Pr(s\ge X)\approx 1/X$ is a random-discriminant model, not a proved distribution for fixed-$p$, prime-order, root-conditioned curves.

\subsection{Isogeny structure}
For $s=1$, $\End(E)=\mathcal{O}_K$ is forced; small rational $\ell$-isogenies are read from roots of $\Phi_\ell(j(E),X)$ modulo $p$. The published P-224 ($s=3$) has four roots at $\ell=3$, locating it on the crater of a height-one 3-volcano. Across all nineteen curves it is not the only nontrivial volcano: prime239 candidates 2320 and 2682 have $s=5$ and only one rational 5-isogeny each. They are on the floors, with $\End(E)=\Z[\pi]$, not $\mathcal{O}_K$; one is published and one skipped. Over $\F_{p^k}$ the Frobenius conductor becomes $sU_k$ and the volcano height is $v_\ell(sU_k)$. The geometric endomorphism ring is unchanged; its conductor fixes the vertex's distance from the crater. Thus not every curve is on the crater over extensions. The tested heights and positions do not yield a separator.

\subsection{Extension-field group orders}
The trial-division study through $k=12$ and the later partial factorizations of $C_3=\#E(\F_{p^3})/n$, $C_4=\#E(\F_{p^4})/\#E(\F_{p^2})=(p-1)^2+t^2$, and $C_6=\#E(\F_{p^6})\,n/(\#E(\F_{p^3})\#E(\F_{p^2}))$ yield no established separating rule. Review found that the old ECM driver could discard later factors in loop mode and conceal process errors. Four extracted composite atoms have been split and all products validated, but the 57 searches have not been rerun; their reported curve budgets are unverified. Prime residuals are now counted among the known primes. For the product of known prime factors below $2^{40}$, the corrected sextic $\log_2$ medians are 51.2 (published) and 21.7 (skipped). An exploratory within-run rank comparison preserving published counts and final endpoints gives a one-sided tail of $11/126\approx8.7\%$, before correction for other statistics. Composite remainders can hide additional factors, and no finite ECM budget guarantees a factor-size cutoff. Moreover $C_6=p^2-(t+1)p+t^2-t+1$ is cheap to compute; $n\nmid C_6$ addresses the MOV condition, not historical non-use.

\subsection{Embedding degree, gcds, symbols}
$\ord_n(p)>20{,}000$ for every viable curve; exact values were obtained for some. In particular, a $2^{222}$ lower bound cannot hold for a 192-bit subgroup. No tested gcd statistic (value, odd part, largest prime factor) ranks published above or below skipped consistently across the four runs. Within prime192 the published curves have $(n/p)=+1,-1,+1$ and $n\equiv1,3,1\pmod 8$; within prime239, $n\equiv3,3,1\pmod8$. No tested symbol or residue separates all four runs.

\subsection{The coefficient $b$ as an integer}
All nineteen primary-model coefficients have complete, checked prime-factor witnesses, with all four alternate P-224 roots retained. The published P-224 coefficient has a 123-bit largest prime factor, not the 149-bit factor of its negative used previously. In the root-2 model, P-224 largest factors are 93, 94, 123 (published), and 203 bits; prime192 still has published values 50, 76, 96 bracketing skipped 77, 95. Four of eight primary published coefficients and six of eleven skipped coefficients are square-free. No tested factor-count, smoothness, or residue statistic explains all labels. The eight published coefficients occupy four classes modulo 7: a specified four-class set has probability $(4/7)^8\approx1.14\%$ under an IID-uniform model, but occupancy of \emph{any} at most four classes has probability $28.4\%$. The classes were selected post hoc, so the smaller figure is not the appropriate significance claim. The character $(b/p)$ is forced positive in the fifteen textbook-root cases; for fixed $p$ it equals the product of the prime-factor characters, with multiplicity, and is not independent information. These finite comparisons do not exhaust numerical criteria.

\subsection{Point structure}
Coordinate comparisons retain root 1 for the unpublished P-224 candidates and use the actual standard coefficient for the published one; those skipped-candidate representatives are explicit models, not recovered historical choices. In this comparison, the smallest $x$ is 0 for sixteen curves, 1 for two skipped P-224 curves, and 3 for the published P-224. Fifteen zeros are forced: the textbook root $b=r^{(p+1)/4}$ is itself a square when $r=-27/c$ is a square. The bitmaps have 118--148 entries set, but an independent fair-bit null ignores this conditioning and algebraic dependencies. The tested bitmaps and small-scalar relations do not give a separator.

\section{Models of the program}\label{sec:program}

\paragraph{RNG alignment.} Published seeds occur at draw offsets $20(i-1)$ from the recovered initialization. Under a shared-stream model with nonoverlapping 20-draw candidate blocks, extra consumption must be a multiple of 20 \emph{cumulatively at those checkpoints}, not at each event. One extra draw followed later by nineteen can restore alignment; intervening candidates need not be aligned. None of the tested fixed-block, event-triggered skip rules fits all nineteen labels. This does not exclude variable-consumption state machines.

\paragraph{Simulated early-abort bugs.} Eleven simple abort rules were tried: the correct abort (control), trace sign flipped, order of the other root, both flipped, $t\equiv0$, $t\equiv\pm1,\pm2$, off-by-one on $n$ in both directions, signed/unsigned confusion, and a twist-order abort. None separates (Appendix~\ref{app:tables}). Bounded versions and the tested single-residue rules likewise fail. These results exclude those models, not every consistent implementation bug or dependence on intermediate arithmetic.

\paragraph{Time limit.} Elkies/Atkin balance for $\ell\le97$: published curves average 9.9 Elkies primes of 24 against 12.5 for skipped. This points against a simple unweighted-count proxy for timeout rejection. Actual historical timings depend on the degrees processed, arithmetic paths, caches, and implementation; no historical runtime measurements were made.

\paragraph{Index-dependent bugs, seed-byte mangling, base-point step, human sorting.} The eleven skipped indices share no residue class, bit, popcount or buffer-offset property disjoint from the eight published. Two skipped seeds (and no published seed) contain a 0x00 byte---consistent with a C-string bug for those two and no others; 0x0A appears on both sides; a bit-length check is excluded by the published P-192 seed (158-bit value). The FreeLIP base-point recipe recovered in \cite{anthropic}, run on every viable prime192/prime239 curve with the published seeds and 600 nearby ones, behaves identically on published and skipped curves (control: the six published seeds reproduce the six published $G_x$). In the three-curve runs the published curves are not the top-3 or bottom-3 under index, seed, $b$, $n$, $t$ or $n-p$.

\section{A union of small bugs? Set cover}\label{sec:cover}

In the original generator-root-1 model, 29 of 135 predicates are false on all eight published curves; the smallest union covering the skips has size 4, with ten covers. There are $\binom{4}{2}\binom{7}{2}=126$ count- and endpoint-preserving relabellings, and 88/126 (69.8\%) admit a cover at least that small. Root 1 is not the published P-224 coefficient. Using root 2 for the four P-224 candidates gives size 4 with 22 covers and 98/126 (77.8\%) doing at least as well. Neither choice establishes the historical skipped-candidate convention. The original count-only Monte Carlo comparison gave 35\%. The adaptive predicate library and root-model choices make these exploratory diagnostics, not confirmatory exclusions of other criteria.

\section{The K-163 generator}\label{sec:k163}

The base points of P-224, P-256, P-384 and P-521 remain unexplained \cite{anthropic}; the six explained generators come from FreeLIP's \texttt{zrandom}. Hashing a phrase directly to $x$ and doubling predicts a short coordinate at $G/2$. The original P-224 scan reconstructed a different point from the published $G_x$ on the opposite-sign model; $nG=O$ did not establish generator identity. The exact standard tuple gives minimum $x$ of 205 bits at $4819G$ and no coordinate at least 32 bits short. A stale, out-of-Hasse P-521 order in the script was also replaced with the standard order; its repeated scan reproduces the cached 502-bit minimum. The other eight generator outputs are retained. All ten declared orders are now checked for primality, the Hasse bound, and scalar multiplication. Their $x$ minima are within two bits of the rough uniform-coordinate minimum.

Over $\F_{2^m}$ the same scan on the ten NIST binary curves ($kG$, $k\le2^{16}$; $G/k$, $k\le2^{12}$; parameters verified by $G\in E$ and $nG=O$) finds nine random-looking curves and one that is not. On K-163 ($y^2+xy=x^3+x^2+1$, cofactor $h=2$),
\[
G/2 = \tfrac{n+1}{2}\,G \quad\text{has}\quad x = \texttt{0x23e21d6019e1211f6bd47ec180256e97},
\]
a 126-bit value---128 bits with two leading zeros---while $y$ is 161 bits. Under a uniform-coordinate heuristic the threshold has probability $2^{-37}$ for a specified $x$. The binary scans contain 69,630 points \emph{per curve}, or 696,300 in total, giving about $5.1\cdot10^{-6}$ expected $x$-hits under that heuristic. Other coordinates, prime-field scans, and post-hoc choices widen the search; this is not a calibrated discovery probability. Choosing a 128-bit $x$, solving for $y$, and multiplying by the cofactor is plausible but unproved.

The original \texttt{binpts.gp} export used \texttt{fromdigits(Vecrev(Vec(f.pol)),2)}: reversing coefficients before interpreting the binary integer. Removing \texttt{Vecrev} gives the standard polynomial-basis coordinate above. Its 126-bit length, computed from the polynomial degree, was unaffected. All ten decimal exports and the hash-search target have been corrected. Earlier negative phrase searches used the wrong target and do not apply to this one.

The subgroup inverse $Q=[h^{-1}\bmod n]G$ is only one possible cofactor-clearing input. All rational preimages are $Q+T$ with $T\in E(\F_{2^m})[h]$. Checking these on all ten curves gives no additional 128-bit-$x$ hit; K-163's other half has 162-bit $x$. The larger subgroup scans do not inspect every such torsion coset.

We have not identified the value or established its novelty. The short half-point is a concrete lead for the separate base-point provenance problem. Anthropic recovered fourteen of fifteen NIST/ANSI seeds, not B-163's, and explained six prime-field generators. A recovered recipe for K-163 would not by itself establish the recipe or date for every binary generator.

\section{Open problems}\label{sec:open}
\begin{enumerate}
\item The eleven skips: seek archival code and printouts, while testing explicit variable-consumption state machines, run-specific criteria, and historical point-counter implementations. Untested numerical rules remain possible.
\item The 128-bit value behind K-163's $G/2$: hash preimage search over 1997-era phrases, encodings and hash functions; or an RNG stream. A match would date and explain the binary-curve base points as the \lrand{} result did the seeds.
\item Base points of P-224, P-256, P-384, P-521, the five B curves and the four other K curves. Extend the search over complete cofactor-preimage sets, field encodings, hash expansions, and RNG layouts rather than only subgroup inverse-multiples.
\item The B-163 seed, which fits no RNG tried in \cite{anthropic}.
\end{enumerate}

\section*{Reproducibility}
Scripts, factor witnesses, order records, and outputs are in the repository. Original point counting used PARI/GP 2.10; review checks and coefficient refinements use PARI/GP 2.19.0. All nineteen viable prime orders have independent primality/Hasse/scalar-multiplication certificates; composite-order scalar checks are consistency checks. The cached factor summarizers run without repeating ECM. Fresh ECM requires GMP-ECM and retains per-attempt transcripts, valid exit statuses, and completed curve counts. Existing extension witnesses retain an explicit legacy-unverified coverage label. The README gives root-independent invocation paths, both P-224 coefficient models, and commands for cached validation, optional msieve replay, fresh ECM, and the corrected point scan.

\section*{Acknowledgements}
\todo{Nadia Heninger for the half-point observation on base points; others.}

\begin{thebibliography}{9}
\bibitem{anthropic} Anthropic Frontier Red Team. Claude reveals ECDSA curve origins. 16 Oct 2026. \url{https://www.anthropic.com/research/ecdsa-curve-origins}
\bibitem{anthropic-crypto} Anthropic Frontier Red Team. Discovering cryptographic weaknesses with Claude. 28 Jul 2026. \url{https://www.anthropic.com/research/discovering-cryptographic-weaknesses}
\bibitem{seeds-solved} F.~Valsorda. The NIST elliptic curves seeds bounty has been claimed. 16 Oct 2026. \url{https://words.filippo.io/seeds-solved/}
\bibitem{bounty} F.~Valsorda. The NIST elliptic curves seeds bounty. \url{https://words.filippo.io/seeds-bounty/}
\bibitem{riddle} N.~Koblitz, A.~Menezes. A riddle wrapped in an enigma. IACR ePrint 2015/1018.
\bibitem{bada55} D.~J. Bernstein, T.~Chou, C.~Chuengsatiansup, A.~H\"ulsing, E.~Lambooij, T.~Lange, R.~Niederhagen, C.~van Vredendaal. How to manipulate curve standards: a white paper for the black hat. IACR ePrint 2014/571.
\bibitem{fips186-2} NIST. FIPS PUB 186-2, Digital Signature Standard. 2000.
\bibitem{x962} ANSI X9.62-1998. Public key cryptography for the financial services industry: the elliptic curve digital signature algorithm (ECDSA).
\bibitem{freelip} A.~K. Lenstra. FreeLIP. \url{https://cr.yp.to/speed/mult/freelip.html}
\end{thebibliography}

\appendix
\section{Tables}\label{app:tables}
\subsection{Class number and $n\pm1$ for all viable curves}
Generic $\log_2 h \approx \log_2\sqrt{4p}$: 129.0 (P-256), 113.0 (P-224), 97.0 (prime192), 120.5 (prime239). $s$ is the conductor of $\Z[\pi]$; $L$ is $L(1,\chi_f)$ by the Euler product to $10^6$; LPF = largest prime factor, in bits.

\begin{center}\small
\begin{tabular}{llrrrrrr}
\toprule
run & cand & $s$ & $L(1,\chi)$ & $\log_2 h$ & LPF$(n{-}1)$ & LPF$(n{+}1)$ & Elkies/24 \\
\midrule
P-256 & 876 & 1 & 0.530 & 126.4 & 242 & 232 & 10 \\
P-256 & \textbf{1842} & 1 & 0.723 & 126.9 & 92 & 215 & 13 \\
P-224 & 777 & 1 & 0.607 & 110.5 & 145 & 182 & 14 \\
P-224 & 1485 & 1 & 0.495 & 110.3 & 74 & 102 & 13 \\
P-224 & 5736 & 1 & 0.540 & 110.4 & 204 & 142 & 11 \\
P-224 & \textbf{7026} & 3 & 0.359 & 108.1 & 196 & 177 & 8 \\
192 & \textbf{139} & 1 & 0.411 & 94.0 & 92 & 108 & 9 \\
192 & 2198 & 1 & 0.498 & 93.6 & 176 & 149 & 11 \\
192 & \textbf{5228} & 1 & 0.352 & 93.8 & 157 & 161 & 6 \\
192 & 5405 & 1 & 0.434 & 94.0 & 132 & 177 & 10 \\
192 & \textbf{6013} & 1 & 0.529 & 93.6 & 64 & 145 & 13 \\
239 & 36 & 1 & 0.436 & 117.1 & 105 & 153 & 12 \\
239 & 1562 & 1 & 0.590 & 118.1 & 219 & 136 & 17 \\
239 & \textbf{2185} & 1 & 0.246 & 116.8 & 69 & 136 & 8 \\
239 & \textbf{2320} & 5 & 0.476 & 115.2 & 128 & 123 & 10 \\
239 & 2491 & 1 & 0.415 & 117.6 & 76 & 131 & 11 \\
239 & 2682 & 5 & 0.558 & 115.7 & 170 & 227 & 17 \\
239 & 2863 & 1 & 0.505 & 117.7 & 183 & 72 & 13 \\
239 & \textbf{3048} & 1 & 0.410 & 117.5 & 110 & 151 & 12 \\
\bottomrule
\end{tabular}
\end{center}
Published curves in bold.

\subsection{Legendre symbols and residues}
\begin{center}\scriptsize
\begin{tabular}{llrrrrrrrrrrr}
\toprule
run & cand & $(n/p)$ & $(p/n)$ & $(t/p)$ & $(t/n)$ & $(b/p)$ & $(-3/n)$ & $(-1/n)$ & $(2/n)$ & $n\,\%\,8$ & $p\,\%\,8$ & $t\,\%\,4$ \\
\midrule
P-256 & 876 & $+$ & $-$ & $+$ & $-$ & $+$ & $+$ & $-$ & $+$ & 7 & 7 & 1 \\
P-256 & \textbf{1842} & $-$ & $-$ & $+$ & $-$ & $+$ & $+$ & $+$ & $+$ & 1 & 7 & 3 \\
P-224 & 777 & $+$ & $+$ & $-$ & $+$ & $+$ & $+$ & $-$ & $-$ & 3 & 1 & 3 \\
P-224 & 1485 & $-$ & $-$ & $+$ & $-$ & $-$ & $+$ & $+$ & $+$ & 1 & 1 & 1 \\
P-224 & 5736 & $-$ & $-$ & $+$ & $-$ & $-$ & $+$ & $+$ & $+$ & 1 & 1 & 1 \\
P-224 & \textbf{7026} & $-$ & $-$ & $+$ & $-$ & $-$ & $+$ & $+$ & $-$ & 5 & 1 & 1 \\
192 & \textbf{139} & $+$ & $+$ & $+$ & $+$ & $+$ & $-$ & $+$ & $+$ & 1 & 7 & 3 \\
192 & 2198 & $+$ & $+$ & $+$ & $+$ & $+$ & $+$ & $+$ & $+$ & 1 & 7 & 3 \\
192 & \textbf{5228} & $-$ & $+$ & $+$ & $-$ & $+$ & $+$ & $-$ & $-$ & 3 & 7 & 1 \\
192 & 5405 & $+$ & $+$ & $-$ & $+$ & $+$ & $+$ & $+$ & $+$ & 1 & 7 & 3 \\
192 & \textbf{6013} & $+$ & $+$ & $+$ & $+$ & $+$ & $+$ & $+$ & $+$ & 1 & 7 & 3 \\
239 & 36 & $+$ & $-$ & $+$ & $-$ & $+$ & $+$ & $-$ & $-$ & 3 & 7 & 1 \\
239 & 1562 & $-$ & $+$ & $-$ & $+$ & $+$ & $+$ & $-$ & $+$ & 7 & 7 & 1 \\
239 & \textbf{2185} & $-$ & $+$ & $-$ & $+$ & $+$ & $+$ & $-$ & $-$ & 3 & 7 & 1 \\
239 & \textbf{2320} & $-$ & $+$ & $+$ & $+$ & $+$ & $+$ & $-$ & $-$ & 3 & 7 & 1 \\
239 & 2491 & $+$ & $+$ & $-$ & $+$ & $+$ & $+$ & $+$ & $-$ & 5 & 7 & 3 \\
239 & 2682 & $+$ & $-$ & $+$ & $-$ & $+$ & $-$ & $-$ & $+$ & 7 & 7 & 1 \\
239 & 2863 & $+$ & $-$ & $+$ & $-$ & $+$ & $-$ & $-$ & $+$ & 7 & 7 & 1 \\
239 & \textbf{3048} & $+$ & $+$ & $+$ & $+$ & $+$ & $-$ & $+$ & $+$ & 1 & 7 & 3 \\
\bottomrule
\end{tabular}
\end{center}

\subsection{Simulated early-abort rules ($\ell\le100$)}
\begin{center}\small
\begin{tabular}{lrr}
\toprule
rule & skipped rejected & published rejected \\
\midrule
(a) $n\equiv0$ --- correct abort (control) & 0/11 & 0/8 \\
(b) $p+1+t\equiv0$ --- trace sign flipped & 9/11 & 7/8 \\
(c) order of the other root of $b$ & 6/11 & 6/8 \\
(d) sign and root both flipped & 3/11 & 1/8 \\
(e) $t\equiv0$ & 5/11 & 4/8 \\
(f) $t\equiv\pm1,\pm2$ & 11/11 & 8/8 \\
(g1) $n\equiv1$ (tests $n-1$) & 11/11 & 8/8 \\
(g2) $n\equiv-1$ (tests $n+1$) & 10/11 & 7/8 \\
(g3) $|t|\equiv0$ & 5/11 & 4/8 \\
(h) twist order $\equiv0$ & 9/11 & 7/8 \\
\bottomrule
\end{tabular}
\end{center}

\subsection{Small multiples of the published generators}
Shortest coordinate among $kG$ ($k\le2^{16}$) and $G/k$ ($k\le2^{16}$ prime-field, $\le2^{12}$ binary), negations included; ``expected'' is the minimum of that many uniform field elements.
\begin{center}\small
\begin{tabular}{lrrrr}
\toprule
curve & $h$ & $x(G/h)$ bits & shortest found (bits) & expected \\
\midrule
P-192 & 1 & -- & 174 & 175 \\ prime192v2 & 1 & -- & 175 & 175 \\ prime192v3 & 1 & -- & 176 & 175 \\
P-224 & 1 & -- & 208 & 207 \\ prime239v1 & 1 & -- & 222 & 222 \\ prime239v2 & 1 & -- & 220 & 222 \\ prime239v3 & 1 & -- & 221 & 222 \\
P-256 & 1 & -- & 237 & 239 \\ P-384 & 1 & -- & 367 & 367 \\ P-521 & 1 & -- & 502 & 504 \\
\midrule
B-163 & 2 & 163 & 147 & 147 \\ \textbf{K-163} & 2 & \textbf{126} & \textbf{126} ($G/2$) & 147 \\
B-233 & 2 & 232 & 215 & 217 \\ K-233 & 4 & full & 219 & 217 \\
B-283 & 2 & 283 & 267 & 267 \\ K-283 & 4 & full & 267 & 267 \\
B-409 & 2 & 408 & 391 & 393 \\ K-409 & 4 & full & 393 & 393 \\
B-571 & 2 & 571 & 554 & 555 \\ K-571 & 4 & full & 555 & 555 \\
\bottomrule
\end{tabular}
\end{center}
For the $h=4$ Koblitz curves $G/4$ lies inside the $G/k$ scan and no coordinate within 32 bits of the field size was found (``full'').


\end{document}
